\subsection{Business Process}

Traditional educational workflows rely heavily on manual 
content creation, fixed assessment schedules, and uniform grading 
criteria, creating systematic bottlenecks in delivering personalised, 
competency-driven feedback at scale. aBridgeAI reengineers this
workflow through a five-stage AI-orchestrated process that separates
programme governance, content preparation, knowledge validation,
practical assessment, and longitudinal reporting into distinct,
traceable phases.

\begin{enumerate}

    \item \textbf{Programme Governance and Enrolment (Admission).}
    Before any content is authored, the structures that give a course
    its meaning must exist. Within a faculty, a Faculty Dean or an
    Educational Manager defines a \textit{learning programme} that
    bundles one or more career pathways --- the ordered, staged
    sequences of courses that constitute a professional competency ---
    and publishes it. Publication freezes a version: subsequent edits
    open a new draft, so a student who enrolled under one version
    continues to see the curriculum they joined. Authorised operators
    then admit students, individually or by bulk import, subject to an
    organisation-configurable ceiling on concurrent enrolments. Each
    programme version nominates one of its pathways as the default,
    and admission places the student on it at once; a student may
    afterwards take on further pathways of their own accord, up to a
    limit the programme sets. Enrolment is what grants access to the
    courses that Stages~2 through~4 operate upon; without it a student
    has no curriculum and the loop cannot begin.

    \item \textbf{Content Ingestion (Preparation).}
    Teachers upload raw learning materials in supported formats 
    including PDF documents, MP4 video recordings, and presentation 
    files. The ingestion pipeline parses each source, extracts typed 
    knowledge units — summaries, concepts, examples, and indexes 
    them in a dense vector store to support semantic retrieval. The 
    resulting structured knowledge base serves as the sole authoritative 
    source for all subsequent AI generation tasks, ensuring that 
    generated assessments are grounded exclusively in 
    teacher-uploaded content.

    \item \textbf{The Learning Loop (Knowledge Validation).}
    Students consume published learning materials and undertake 
    AI-generated quizzes to validate baseline knowledge retention. 
    Each quiz card carries a teacher-confirmed expected response time 
    $T_{\text{exp}}$. The system derives the quality response value 
    $Q$ algorithmically from response correctness, hint usage, and 
    the time ratio $\rho = T_{\text{actual}} / T_{\text{exp}}$, 
    updating the student's per-card Easiness Factor $EF$ via the 
    hybrid SM-2 scheduling model. Students remain in the learning 
    loop — receiving SR reminders at algorithmically 
    determined intervals — until their $EF$ distribution across the 
    lesson's cards satisfies the teacher-configured unlock condition.

    \item \textbf{The Application Loop (Practical Assessment).}
    Upon satisfying the quiz unlock condition, students gain access 
    to the lesson's interview section. Interview questions are drawn 
    from a pre-generated, teacher-reviewed question bank aligned to 
    teacher-defined outcome requirements. Students engage with the 
    interview in text or voice mode, with each session conducted 
    independently under a dedicated model instance. Upon session 
    completion, the full transcript is submitted to the AI assessment 
    model in a single batch call, which evaluates whether each defined 
    outcome was satisfied. The student receives a binary pass or fail 
    verdict; no rubric, per-outcome breakdown, or model reasoning is 
    exposed. A lesson is considered complete only when both the quiz 
    unlock condition and the interview pass condition are satisfied.

    \item \textbf{Evaluation and Reporting (Feedback).}
    The system continuously aggregates per-card SM-2 state into three 
    learner-facing metrics: the Knowledge Retention Estimate 
    $\hat{R}_{s,\ell}$, Progression Readiness, and Review Compliance 
    Rate $\rho_{s,\ell}$. These metrics are surfaced in per-student 
    and class-wide dashboard views accessible to both students and 
    teachers. At the career level, course completion percentages are 
    aggregated into a Career Readiness score reported to the student,
    the Educational Manager, and the Faculty Dean who owns the
    programme. Teachers are additionally notified
    of cohort-level patterns — including systematic retention gaps and
    declining compliance rates — enabling intervention at the content
    or configuration level rather than purely at the individual
    student level.

    Reporting also feeds a slower, governance-level response. A
    student whose progression shows that they are on the wrong pathway
    may request a change of career pathway, or --- where they carry
    more than one --- that a pathway be dropped, in each case with a
    written justification; the Faculty Dean who owns the programme
    acknowledges the request and then approves or rejects it against
    the same progression evidence, and may never adjudicate their own.
    Taking on a pathway is thus the student's to decide and laying one
    down is not, and no approval may leave a student with none. An
    approval snapshots the progress earned on the outgoing pathway
    and preserves completed courses as awards that carry across, so
    that work shared between two pathways is credited once and
    never repeated. The decision, and the reason behind it, are returned to
    the student.

\end{enumerate}

The five stages form two nested feedback loops that turn at different
speeds. The \textit{pedagogical loop} closes within a course: content
ingestion informs quiz generation, quiz performance governs interview
access, interview outcomes determine curriculum completion, and
reporting surfaces the gaps that guide both student remediation and
teacher reconfiguration of content and assessment criteria. It turns
with every review session.

The \textit{governance loop} closes around the programme and turns
rarely. Enrolment in Stage~1 fixes the pathway a student is measured
against; the reporting of Stage~5 evidences whether that pathway
remains the right one; and a dean-adjudicated pathway change feeds
back into Stage~1 by opening a new attempt under a different pathway,
without discarding what the student has already earned. Separating the
two matters because they fail differently: a defect in the
pedagogical loop misjudges a single student's readiness on one lesson,
whereas a defect in the governance loop misdirects an entire
programme of study.

